\section{Proof of the main theorem}
\label{sec:conclusion}

\begin{proof}[Proof of Theorem~\ref{thm:main}]
Take the algebraic rank-two bundle $G$ of
Proposition~\ref{prop:construction}. By
Proposition~\ref{prop:ampleness}, every $E_m=f_m^*G\otimes P$ is ample
in the quotient convention.

If the set of positive integers $m$ for which $E_m$ has a smooth
strictly Griffiths-positive metric were infinite, it would contain
a sequence tending to infinity. Proposition~\ref{prop:limit} would
then give a finite, nonnegative, nonzero psh function on the total
space of $G^\vee$ that is a Hermitian positive semidefinite quadratic
form on every fiber. Proposition~\ref{prop:obstruction} excludes
precisely such a function. Consequently that set of integers is
finite. Choose $m_0$ larger than all its elements, taking $m_0=1$ if
the set is empty. Every $m\ge m_0$ gives the required counterexample.
\end{proof}

The limit in the proof is permitted to lose rank, but not to vanish
identically or to lose its quadratic structure at any base point.
These are exactly the properties needed to transfer the
graph relations to a nonzero invariant measure. The argument requires no uniform estimates for the input metrics or
their curvature, and it does not determine an effective value of $m_0$.
